9.2 Topological Spaces
\hfill 13
We shall be dealing mainly with separable spaces.
We now remark that, since the definition of a neighborhood of a point in a topo-
logical space is verbally the same as the definition of a neighborhood of a point in
a metric space, the concepts of interior point, exterior point, boundary point, and
limit point of a set, and the concept of the closure of a set, all of which use only the
concept of a neighborhood, can be carried over without any changes to the case of
an arbitrary topological space.
Moreover, as can be seen from the proof of Proposition 2 in Sect. 7.1, it is also
true that a set in a Hausdorff space is closed if and only if it contains all its limit
points.

\noindent
\textbf{9.2.2 Subspaces of a Topological Space}

Let $(X, \tau_X)$ be a topological space and $Y$ a subset of $X$. The topology $\tau_X$ makes
it possible to define the following topology $\tau_Y$ in $Y$, called the induced or relative
topology on $Y \subset X$.
We define an open set in $Y$ to be any set $G_Y$ of the form $G_Y = Y \cap G_X$, where
$G_X$ is an open set in $X$.
It is not difficult to verify that the system $\tau_Y$ of subsets of $Y$ that arises in this
way satisfies the axioms for open sets in a topological space.
As one can see, the definition of open sets $G_Y$ in $Y$ agrees with the one we
obtained in Sect. 9.1.3 for the case when $Y$ is a subspace of a metric space $X$.

\noindent
\textbf{Definition 9} A subset $Y \subset X$ of a topological space $(X, \tau)$ with the topology $\tau_Y$
induced on $Y$ is called a subspace of the topological space $X$.
It is clear that a set that is open in $(Y, \tau_Y)$ is not necessarily open in $(X, \tau_X)$.

\noindent
\textbf{9.2.3 The Direct Product of Topological Spaces}

If $(X_1, \tau_1)$ and $(X_2, \tau_2)$ are two topological spaces with systems of open sets $\tau_1 =
\{G_1\}$ and $\tau_2 = \{G_2\}$, we can introduce a topology on $X_1 \times X_2$ by taking as the base
the sets of the form $G_1 \times G_2$.

\noindent
\textbf{Definition 10} The topological space $(X_1 \times X_2, \tau_1 \times \tau_2)$ whose topology has the
base consisting of sets of the form $G_1 \times G_2$, where $G_i$ is an open set in the topo-
logical space $(X_i, \tau_i)$, $i = 1, 2$, is called the direct product of the topological spaces
$(X_1, \tau_1)$ and $(X_2, \tau_2)$.

\noindent
\textbf{Example 9} If $\R = \R^1$ and $\R^2$ are considered with their standard topologies, then,
as one can see, $\R^2$ is the direct product $\R^1 \times \R^1$. For every open set in $\R^2$ can be
obtained, for example, as the union of "square" neighborhoods of all its points. And